\chapitre{Fenêtres homologues du chromosome humain}
Les fenêtres homologues comparent des segments correspondants de GRCh38.p14 et T2T-CHM13v2.0~\cite{t2t,chmreleases}. Sur le chromosome 21, les intervalles commencent en 34\,000\,000 et 32\,381\,795, en coordonnées de base zéro. Les ancres de 41 bases qui encadrent chaque intervalle sont uniques sur leur chromosome. La chaîne de correspondance UCSC place leurs extrémités dans le même ordre sur le brin positif. Les insertions et les délétions internes restent dans les fenêtres.
\chart{11_chromosome21.png}{Les fenêtres homologues et la copie du chromosome portant des insertions connues comparent les durées et les empreintes des méthodes.}{fig:humain}{bench/human.sbatch}
Sur la figure~\ref{fig:humain}, deux processus par méthode donnent les temps médians sur Grace et GH200. Pour les fenêtres de deux puis quatre millions de bases de référence, \(d=5\,599\) puis \(12\,482\). La portée \(R\) de l'index atteint 290 puis 622. Les deux calculs NvPD emploient donc un lancement par ligne. Leur temps passe de 18,350 à 90,891\,s, allocations, transferts et initialisation CUDA compris. Sur un cœur Grace, WFA2 demande 0,054 puis 0,348\,s et edlib 0,790 puis 2,991\,s.

Le contrôle entier ajoute 64 bases à une copie des 45\,090\,682 bases de CHM13. Le chemin construit coûte 64 et la différence de longueur impose \(d\ge64\). WFA2 obtient ce score en 11,34\,ms de calcul. Le processus prend 0,476\,s avec la lecture des deux chaînes.
\synthese{Sur les fenêtres homologues, un quadruplement approximatif du nombre de cases multiplie le temps du GPU par 4,95.}{À l'échelle d'un chromosome, la divergence entre versions d'assemblage peut varier de plusieurs ordres de grandeur.}

\chapitre{Comparaison des assemblages chromosomiques entiers}
Les versions v1.0 et v2.0 de CHM13 permettent de mesurer le coût d'une révision d'assemblage sur le chromosome 20. Les entrées \texttt{CP068258.1} et \texttt{CP068258.2} comptent 66\,210\,247 et 66\,210\,255 bases.\footnote{Voir l'\href{https://github.com/marbl/CHM13/blob/master/Earlier_assembly_releases_and_associated_data.md}{historique des assemblages CHM13 du consortium T2T}.} Le chromosome 20 porte déjà l'accession \texttt{CP068258.2} dans la version intermédiaire 1.1. Sur le chromosome 21, GRCh38 et CHM13 représentent deux références distinctes et comportent des bases ambiguës traitées comme des caractères ordinaires.

Une borne exploite l'histogramme des lettres. Pour \(u_c=\#_A(c)-\#_B(c)\) et \(x_+=\max(0,x)\), chaque édition efface au plus une unité d'excédent et une de déficit. Ainsi,
\[
 d(A,B)\ge\max\!\left(\sum_c(u_c)_+,\ \sum_c(-u_c)_+\right).
\]
Une borne supérieure suit un alignement global réalisable. Sur chaque bloc de la chaîne UCSC 23, il paie les désaccords. Dans chaque intervalle entre blocs, il choisit le moins coûteux des deux alignements ancrés à gauche ou à droite, puis paie l'écart de longueur. Les extrémités suivent la même règle. La somme des coûts des blocs et des intervalles donne le coût du chemin entier.
\chart{18_chromosomes_entiers.png}{Les durées sur le chromosome 20 et les bornes de distance sur le chromosome 21 comparent deux niveaux de divergence entre assemblages.}{fig:chrentiers}{bench/whole_chromosomes.py}
Les trois appels par méthode sur un cœur EPYC donnent \(d=30\) sur le chromosome 20. Le temps médian de calcul vaut 1,280\,s pour edlib et 0,092\,s pour WFA2. Avec lecture du FASTA, les processus prennent 2,679 et 1,484\,s. Sur le chromosome 21 de la figure~\ref{fig:chrentiers}, chaque processus atteint la limite murale de 120\,s. L'écart de longueur seul vaut 1\,619\,301 bases.
\synthese{Les fréquences des bases et le chemin construit encadrent la distance du chromosome 21 entre 6\,621\,364 et 16\,015\,052 éditions.}{Sur cette paire, la composition des séquences relève la borne de longueur d'un facteur 4,09.}


\chapitre{Géométrie des distances entre génomes de phage}
La reconstruction mesure dans quelle proportion les chemins d'un arbre reproduisent les distances globales de 64 séquences de phage \(\phi\)X174. Une origine différente sur le génome circulaire peut modifier la distance entre les chaînes déposées. Une même édition observée peut aussi résulter de plusieurs événements évolutifs.

La méthode UPGMA~\cite{upgma} joint les groupes \(U,V\) de distance moyenne minimale et place leur parent à la hauteur \(d(U,V)/2\). Pour un troisième groupe \(W\), la nouvelle distance moyenne est pondérée par les effectifs. Neighbor-Joining joint la paire qui minimise \(Q_{ij}\), où \(r\) est le nombre de groupes actifs et \(s_i=\sum_k d_{ik}\)~\cite{nj}.
\[
 d(U\cup V,W)=\frac{|U|d(U,W)+|V|d(V,W)}{|U|+|V|},\qquad
 Q_{ij}=(r-2)d_{ij}-s_i-s_j.
\]
Dans NJ, la première branche mesure \(d_{ij}/2+(s_i-s_j)/(2(r-2))\), la seconde complète \(d_{ij}\). À chaque fusion, les deux implémentations examinent toutes les paires actives. Pour \(N\) feuilles, elles utilisent \(O(N^3)\) opérations et \(O(N^2)\) distances stockées. UPGMA impose des feuilles équidistantes de la racine. Sur trois feuilles, les deux plus grandes distances doivent alors coïncider. L'additivité d'un arbre impose sur quatre feuilles l'égalité des deux plus grandes sommes parmi \(d_{ij}+d_{kl},d_{ik}+d_{jl},d_{il}+d_{jk}\).
\chart{16_phylogenie.png}{Les écarts aux conditions d'ultramétrie et d'additivité sont comparés aux résidus entre distances d'édition et longueurs des chemins.}{fig:phage}{scripts/phylo_fit.py}
La figure~\ref{fig:phage} porte sur les \(C(64,3)=41\,664\) triplets et \(C(64,4)=635\,376\) quartets. Les conditions sont satisfaites pour 1\,661 triplets et 245\,847 quartets. La racine de l'écart quadratique moyen vaut \(\sqrt{P^{-1}\sum_{i<j}(d_{ij}-t_{ij})^2}\), où \(t_{ij}\) est la longueur du chemin entre feuilles.
\synthese{Sur les 2\,016 paires, la racine de l'écart quadratique moyen vaut 14,89 éditions pour UPGMA et 3,56 pour NJ.}{La représentation des branches permet de localiser ces écarts dans les deux topologies.}

\chapitre{Topologie des arbres de phage}
Les phylogrammes de la figure~\ref{fig:arbres} utilisent la longueur cumulée des branches comme abscisse. Les échelles horizontales sont propres à chaque panneau. La longueur patristique se calcule par \(t_{ij}=\sum_{e\in\operatorname{chemin}(i,j)}\ell_e\). Le choix du point de dessin initial de NJ laisse ces sommes invariantes.
\begin{figure}[H]\centering
\resizebox{!}{.68\textheight}{\begin{tikzpicture}[x=1mm,y=-1mm,line width=.3pt]
\node[anchor=west,font=\small\bfseries,inner sep=0pt] at (0.000,-6) {a) UPGMA};
\draw[draw=powderBlue!50!black] (5.447,0.000) -- (34.880,0.000);
\draw[draw=powderBlue!50!black] (6.097,3.240) -- (34.880,3.240);
\draw[draw=powderBlue!50!black] (6.543,6.480) -- (34.880,6.480);
\draw[draw=powderBlue!50!black] (15.886,9.720) -- (34.880,9.720);
\draw[draw=powderBlue!50!black] (15.886,12.960) -- (34.880,12.960);
\draw[draw=powderBlue!50!black] (15.886,9.720) -- (15.886,12.960);
\draw[draw=powderBlue!50!black] (7.494,11.340) -- (15.886,11.340);
\draw[draw=powderBlue!50!black] (20.257,16.200) -- (34.880,16.200);
\draw[draw=powderBlue!50!black] (24.704,19.440) -- (34.880,19.440);
\draw[draw=powderBlue!50!black] (26.287,22.680) -- (34.880,22.680);
\draw[draw=powderBlue!50!black] (26.287,25.920) -- (34.880,25.920);
\draw[draw=powderBlue!50!black] (26.287,22.680) -- (26.287,25.920);
\draw[draw=powderBlue!50!black] (24.704,24.300) -- (26.287,24.300);
\draw[draw=powderBlue!50!black] (24.704,19.440) -- (24.704,24.300);
\draw[draw=powderBlue!50!black] (20.257,21.870) -- (24.704,21.870);
\draw[draw=powderBlue!50!black] (20.257,16.200) -- (20.257,21.870);
\draw[draw=powderBlue!50!black] (8.570,19.035) -- (20.257,19.035);
\draw[draw=powderBlue!50!black] (8.752,29.160) -- (34.880,29.160);
\draw[draw=powderBlue!50!black] (9.265,32.400) -- (34.880,32.400);
\draw[draw=powderBlue!50!black] (11.851,35.640) -- (34.880,35.640);
\draw[draw=powderBlue!50!black] (12.207,38.880) -- (34.880,38.880);
\draw[draw=powderBlue!50!black] (12.984,42.120) -- (34.880,42.120);
\draw[draw=powderBlue!50!black] (13.121,45.360) -- (34.880,45.360);
\draw[draw=powderBlue!50!black] (14.356,48.600) -- (34.880,48.600);
\draw[draw=powderBlue!50!black] (14.502,51.840) -- (34.880,51.840);
\draw[draw=powderBlue!50!black] (15.674,55.080) -- (34.880,55.080);
\draw[draw=powderBlue!50!black] (15.900,58.320) -- (34.880,58.320);
\draw[draw=powderBlue!50!black] (15.961,61.560) -- (34.880,61.560);
\draw[draw=powderBlue!50!black] (16.151,64.800) -- (34.880,64.800);
\draw[draw=powderBlue!50!black] (16.903,68.040) -- (34.880,68.040);
\draw[draw=powderBlue!50!black] (18.247,71.280) -- (34.880,71.280);
\draw[draw=powderBlue!50!black] (20.408,74.520) -- (34.880,74.520);
\draw[draw=powderBlue!50!black] (27.945,77.760) -- (34.880,77.760);
\draw[draw=powderBlue!50!black] (32.618,81.000) -- (34.880,81.000);
\draw[draw=powderBlue!50!black] (33.071,84.240) -- (34.880,84.240);
\draw[draw=powderBlue!50!black] (33.071,87.480) -- (34.880,87.480);
\draw[draw=powderBlue!50!black] (33.071,84.240) -- (33.071,87.480);
\draw[draw=powderBlue!50!black] (32.618,85.860) -- (33.071,85.860);
\draw[draw=powderBlue!50!black] (32.618,81.000) -- (32.618,85.860);
\draw[draw=powderBlue!50!black] (27.945,83.430) -- (32.618,83.430);
\draw[draw=powderBlue!50!black] (27.945,77.760) -- (27.945,83.430);
\draw[draw=powderBlue!50!black] (22.001,80.595) -- (27.945,80.595);
\draw[draw=powderBlue!50!black] (27.644,90.720) -- (34.880,90.720);
\draw[draw=powderBlue!50!black] (33.071,93.960) -- (34.880,93.960);
\draw[draw=powderBlue!50!black] (33.071,97.200) -- (34.880,97.200);
\draw[draw=powderBlue!50!black] (33.071,93.960) -- (33.071,97.200);
\draw[draw=powderBlue!50!black] (27.644,95.580) -- (33.071,95.580);
\draw[draw=powderBlue!50!black] (27.644,90.720) -- (27.644,95.580);
\draw[draw=powderBlue!50!black] (23.613,93.150) -- (27.644,93.150);
\draw[draw=powderBlue!50!black] (30.583,100.440) -- (34.880,100.440);
\draw[draw=powderBlue!50!black] (33.523,103.680) -- (34.880,103.680);
\draw[draw=powderBlue!50!black] (33.523,106.920) -- (34.880,106.920);
\draw[draw=powderBlue!50!black] (33.523,103.680) -- (33.523,106.920);
\draw[draw=powderBlue!50!black] (30.583,105.300) -- (33.523,105.300);
\draw[draw=powderBlue!50!black] (30.583,100.440) -- (30.583,105.300);
\draw[draw=powderBlue!50!black] (26.758,102.870) -- (30.583,102.870);
\draw[draw=powderBlue!50!black] (29.212,110.160) -- (34.880,110.160);
\draw[draw=powderBlue!50!black] (32.618,113.400) -- (34.880,113.400);
\draw[draw=powderBlue!50!black] (32.618,116.640) -- (34.880,116.640);
\draw[draw=powderBlue!50!black] (32.618,113.400) -- (32.618,116.640);
\draw[draw=powderBlue!50!black] (30.862,115.020) -- (32.618,115.020);
\draw[draw=powderBlue!50!black] (31.488,119.880) -- (34.880,119.880);
\draw[draw=powderBlue!50!black] (33.674,123.120) -- (34.880,123.120);
\draw[draw=powderBlue!50!black] (33.975,126.360) -- (34.880,126.360);
\draw[draw=powderBlue!50!black] (33.975,129.600) -- (34.880,129.600);
\draw[draw=powderBlue!50!black] (33.975,132.840) -- (34.880,132.840);
\draw[draw=powderBlue!50!black] (33.975,129.600) -- (33.975,132.840);
\draw[draw=powderBlue!50!black] (33.975,131.220) -- (33.975,131.220);
\draw[draw=powderBlue!50!black] (33.975,126.360) -- (33.975,131.220);
\draw[draw=powderBlue!50!black] (33.674,128.790) -- (33.975,128.790);
\draw[draw=powderBlue!50!black] (33.674,123.120) -- (33.674,128.790);
\draw[draw=powderBlue!50!black] (32.505,125.955) -- (33.674,125.955);
\draw[draw=powderBlue!50!black] (33.975,136.080) -- (34.880,136.080);
\draw[draw=powderBlue!50!black] (33.975,139.320) -- (34.880,139.320);
\draw[draw=powderBlue!50!black] (33.975,136.080) -- (33.975,139.320);
\draw[draw=powderBlue!50!black] (33.221,137.700) -- (33.975,137.700);
\draw[draw=powderBlue!50!black] (33.613,142.560) -- (34.880,142.560);
\draw[draw=powderBlue!50!black] (34.880,145.800) -- (34.880,145.800);
\draw[draw=powderBlue!50!black] (34.880,149.040) -- (34.880,149.040);
\draw[draw=powderBlue!50!black] (34.880,145.800) -- (34.880,149.040);
\draw[draw=powderBlue!50!black] (34.126,147.420) -- (34.880,147.420);
\draw[draw=powderBlue!50!black] (34.201,152.280) -- (34.880,152.280);
\draw[draw=powderBlue!50!black] (34.427,155.520) -- (34.880,155.520);
\draw[draw=powderBlue!50!black] (34.427,158.760) -- (34.880,158.760);
\draw[draw=powderBlue!50!black] (34.427,155.520) -- (34.427,158.760);
\draw[draw=powderBlue!50!black] (34.201,157.140) -- (34.427,157.140);
\draw[draw=powderBlue!50!black] (34.201,152.280) -- (34.201,157.140);
\draw[draw=powderBlue!50!black] (34.126,154.710) -- (34.201,154.710);
\draw[draw=powderBlue!50!black] (34.126,147.420) -- (34.126,154.710);
\draw[draw=powderBlue!50!black] (33.613,151.065) -- (34.126,151.065);
\draw[draw=powderBlue!50!black] (33.613,142.560) -- (33.613,151.065);
\draw[draw=powderBlue!50!black] (33.221,146.812) -- (33.613,146.812);
\draw[draw=powderBlue!50!black] (33.221,137.700) -- (33.221,146.812);
\draw[draw=powderBlue!50!black] (32.505,142.256) -- (33.221,142.256);
\draw[draw=powderBlue!50!black] (32.505,125.955) -- (32.505,142.256);
\draw[draw=powderBlue!50!black] (31.488,134.106) -- (32.505,134.106);
\draw[draw=powderBlue!50!black] (31.488,119.880) -- (31.488,134.106);
\draw[draw=powderBlue!50!black] (30.862,126.993) -- (31.488,126.993);
\draw[draw=powderBlue!50!black] (30.862,115.020) -- (30.862,126.993);
\draw[draw=powderBlue!50!black] (29.212,121.006) -- (30.862,121.006);
\draw[draw=powderBlue!50!black] (29.212,110.160) -- (29.212,121.006);
\draw[draw=powderBlue!50!black] (26.758,115.583) -- (29.212,115.583);
\draw[draw=powderBlue!50!black] (26.758,102.870) -- (26.758,115.583);
\draw[draw=powderBlue!50!black] (23.613,109.227) -- (26.758,109.227);
\draw[draw=powderBlue!50!black] (23.613,93.150) -- (23.613,109.227);
\draw[draw=powderBlue!50!black] (22.001,101.188) -- (23.613,101.188);
\draw[draw=powderBlue!50!black] (22.001,80.595) -- (22.001,101.188);
\draw[draw=powderBlue!50!black] (20.408,90.892) -- (22.001,90.892);
\draw[draw=powderBlue!50!black] (20.408,74.520) -- (20.408,90.892);
\draw[draw=powderBlue!50!black] (18.247,82.706) -- (20.408,82.706);
\draw[draw=powderBlue!50!black] (18.247,71.280) -- (18.247,82.706);
\draw[draw=powderBlue!50!black] (16.903,76.993) -- (18.247,76.993);
\draw[draw=powderBlue!50!black] (16.903,68.040) -- (16.903,76.993);
\draw[draw=powderBlue!50!black] (16.151,72.516) -- (16.903,72.516);
\draw[draw=powderBlue!50!black] (16.151,64.800) -- (16.151,72.516);
\draw[draw=powderBlue!50!black] (15.961,68.658) -- (16.151,68.658);
\draw[draw=powderBlue!50!black] (15.961,61.560) -- (15.961,68.658);
\draw[draw=powderBlue!50!black] (15.900,65.109) -- (15.961,65.109);
\draw[draw=powderBlue!50!black] (15.900,58.320) -- (15.900,65.109);
\draw[draw=powderBlue!50!black] (15.674,61.715) -- (15.900,61.715);
\draw[draw=powderBlue!50!black] (15.674,55.080) -- (15.674,61.715);
\draw[draw=powderBlue!50!black] (14.502,58.397) -- (15.674,58.397);
\draw[draw=powderBlue!50!black] (14.502,51.840) -- (14.502,58.397);
\draw[draw=powderBlue!50!black] (14.356,55.119) -- (14.502,55.119);
\draw[draw=powderBlue!50!black] (14.356,48.600) -- (14.356,55.119);
\draw[draw=powderBlue!50!black] (13.121,51.859) -- (14.356,51.859);
\draw[draw=powderBlue!50!black] (13.121,45.360) -- (13.121,51.859);
\draw[draw=powderBlue!50!black] (12.984,48.610) -- (13.121,48.610);
\draw[draw=powderBlue!50!black] (12.984,42.120) -- (12.984,48.610);
\draw[draw=powderBlue!50!black] (12.207,45.365) -- (12.984,45.365);
\draw[draw=powderBlue!50!black] (12.207,38.880) -- (12.207,45.365);
\draw[draw=powderBlue!50!black] (11.851,42.122) -- (12.207,42.122);
\draw[draw=powderBlue!50!black] (11.851,35.640) -- (11.851,42.122);
\draw[draw=powderBlue!50!black] (9.265,38.881) -- (11.851,38.881);
\draw[draw=powderBlue!50!black] (9.265,32.400) -- (9.265,38.881);
\draw[draw=powderBlue!50!black] (8.752,35.641) -- (9.265,35.641);
\draw[draw=powderBlue!50!black] (8.752,29.160) -- (8.752,35.641);
\draw[draw=powderBlue!50!black] (8.570,32.400) -- (8.752,32.400);
\draw[draw=powderBlue!50!black] (8.570,19.035) -- (8.570,32.400);
\draw[draw=powderBlue!50!black] (7.494,25.718) -- (8.570,25.718);
\draw[draw=powderBlue!50!black] (7.494,11.340) -- (7.494,25.718);
\draw[draw=powderBlue!50!black] (6.543,18.529) -- (7.494,18.529);
\draw[draw=powderBlue!50!black] (6.543,6.480) -- (6.543,18.529);
\draw[draw=powderBlue!50!black] (6.097,12.504) -- (6.543,12.504);
\draw[draw=powderBlue!50!black] (6.097,3.240) -- (6.097,12.504);
\draw[draw=powderBlue!50!black] (5.447,7.872) -- (6.097,7.872);
\draw[draw=powderBlue!50!black] (5.447,0.000) -- (5.447,7.872);
\draw[draw=powderBlue!50!black] (0.000,3.936) -- (5.447,3.936);
\draw[draw=powderBlue!50!black] (16.338,162.000) -- (34.880,162.000);
\draw[draw=powderBlue!50!black] (31.714,165.240) -- (34.880,165.240);
\draw[draw=powderBlue!50!black] (31.714,168.480) -- (34.880,168.480);
\draw[draw=powderBlue!50!black] (31.714,165.240) -- (31.714,168.480);
\draw[draw=powderBlue!50!black] (16.338,166.860) -- (31.714,166.860);
\draw[draw=powderBlue!50!black] (16.338,162.000) -- (16.338,166.860);
\draw[draw=powderBlue!50!black] (0.976,164.430) -- (16.338,164.430);
\draw[draw=powderBlue!50!black] (24.704,171.720) -- (34.880,171.720);
\draw[draw=powderBlue!50!black] (31.714,174.960) -- (34.880,174.960);
\draw[draw=powderBlue!50!black] (31.714,178.200) -- (34.880,178.200);
\draw[draw=powderBlue!50!black] (31.714,174.960) -- (31.714,178.200);
\draw[draw=powderBlue!50!black] (24.704,176.580) -- (31.714,176.580);
\draw[draw=powderBlue!50!black] (24.704,171.720) -- (24.704,176.580);
\draw[draw=powderBlue!50!black] (11.797,174.150) -- (24.704,174.150);
\draw[draw=powderBlue!50!black] (31.262,181.440) -- (34.880,181.440);
\draw[draw=powderBlue!50!black] (32.166,184.680) -- (34.880,184.680);
\draw[draw=powderBlue!50!black] (32.166,187.920) -- (34.880,187.920);
\draw[draw=powderBlue!50!black] (32.166,184.680) -- (32.166,187.920);
\draw[draw=powderBlue!50!black] (31.262,186.300) -- (32.166,186.300);
\draw[draw=powderBlue!50!black] (31.262,181.440) -- (31.262,186.300);
\draw[draw=powderBlue!50!black] (16.278,183.870) -- (31.262,183.870);
\draw[draw=powderBlue!50!black] (22.217,191.160) -- (34.880,191.160);
\draw[draw=powderBlue!50!black] (28.850,194.400) -- (34.880,194.400);
\draw[draw=powderBlue!50!black] (29.679,197.640) -- (34.880,197.640);
\draw[draw=powderBlue!50!black] (31.714,200.880) -- (34.880,200.880);
\draw[draw=powderBlue!50!black] (31.714,204.120) -- (34.880,204.120);
\draw[draw=powderBlue!50!black] (31.714,200.880) -- (31.714,204.120);
\draw[draw=powderBlue!50!black] (29.679,202.500) -- (31.714,202.500);
\draw[draw=powderBlue!50!black] (29.679,197.640) -- (29.679,202.500);
\draw[draw=powderBlue!50!black] (28.850,200.070) -- (29.679,200.070);
\draw[draw=powderBlue!50!black] (28.850,194.400) -- (28.850,200.070);
\draw[draw=powderBlue!50!black] (22.217,197.235) -- (28.850,197.235);
\draw[draw=powderBlue!50!black] (22.217,191.160) -- (22.217,197.235);
\draw[draw=powderBlue!50!black] (16.278,194.198) -- (22.217,194.198);
\draw[draw=powderBlue!50!black] (16.278,183.870) -- (16.278,194.198);
\draw[draw=powderBlue!50!black] (11.797,189.034) -- (16.278,189.034);
\draw[draw=powderBlue!50!black] (11.797,174.150) -- (11.797,189.034);
\draw[draw=powderBlue!50!black] (0.976,181.592) -- (11.797,181.592);
\draw[draw=powderBlue!50!black] (0.976,164.430) -- (0.976,181.592);
\draw[draw=powderBlue!50!black] (0.000,173.011) -- (0.976,173.011);
\draw[draw=powderBlue!50!black] (0.000,3.936) -- (0.000,173.011);
\draw[densely dotted,black!36] (34.880,0.000) -- (35.880,0.000);
\node[anchor=west,inner sep=0pt,font=\rmfamily\fontsize{7}{7}\selectfont] at (36.880,0.000) {EF380022.1};
\draw[densely dotted,black!36] (34.880,3.240) -- (35.880,3.240);
\node[anchor=west,inner sep=0pt,font=\rmfamily\fontsize{7}{7}\selectfont] at (36.880,3.240) {EF380026.1};
\draw[densely dotted,black!36] (34.880,6.480) -- (35.880,6.480);
\node[anchor=west,inner sep=0pt,font=\rmfamily\fontsize{7}{7}\selectfont] at (36.880,6.480) {EF380023.1};
\draw[densely dotted,black!36] (34.880,9.720) -- (35.880,9.720);
\node[anchor=west,inner sep=0pt,font=\rmfamily\fontsize{7}{7}\selectfont] at (36.880,9.720) {EF380024.1};
\draw[densely dotted,black!36] (34.880,12.960) -- (35.880,12.960);
\node[anchor=west,inner sep=0pt,font=\rmfamily\fontsize{7}{7}\selectfont] at (36.880,12.960) {EF380025.1};
\draw[densely dotted,black!36] (34.880,16.200) -- (35.880,16.200);
\node[anchor=west,inner sep=0pt,font=\rmfamily\fontsize{7}{7}\selectfont] at (36.880,16.200) {EF380029.1};
\draw[densely dotted,black!36] (34.880,19.440) -- (35.880,19.440);
\node[anchor=west,inner sep=0pt,font=\rmfamily\fontsize{7}{7}\selectfont] at (36.880,19.440) {EF380021.1};
\draw[densely dotted,black!36] (34.880,22.680) -- (35.880,22.680);
\node[anchor=west,inner sep=0pt,font=\rmfamily\fontsize{7}{7}\selectfont] at (36.880,22.680) {EF380028.1};
\draw[densely dotted,black!36] (34.880,25.920) -- (35.880,25.920);
\node[anchor=west,inner sep=0pt,font=\rmfamily\fontsize{7}{7}\selectfont] at (36.880,25.920) {EF380030.1};
\draw[densely dotted,black!36] (34.880,29.160) -- (35.880,29.160);
\node[anchor=west,inner sep=0pt,font=\rmfamily\fontsize{7}{7}\selectfont] at (36.880,29.160) {EF380027.1};
\draw[densely dotted,black!36] (34.880,32.400) -- (35.880,32.400);
\node[anchor=west,inner sep=0pt,font=\rmfamily\fontsize{7}{7}\selectfont] at (36.880,32.400) {EF380012.1};
\draw[densely dotted,black!36] (34.880,35.640) -- (35.880,35.640);
\node[anchor=west,inner sep=0pt,font=\rmfamily\fontsize{7}{7}\selectfont] at (36.880,35.640) {EF380013.1};
\draw[densely dotted,black!36] (34.880,38.880) -- (35.880,38.880);
\node[anchor=west,inner sep=0pt,font=\rmfamily\fontsize{7}{7}\selectfont] at (36.880,38.880) {EF380011.1};
\draw[densely dotted,black!36] (34.880,42.120) -- (35.880,42.120);
\node[anchor=west,inner sep=0pt,font=\rmfamily\fontsize{7}{7}\selectfont] at (36.880,42.120) {EF380014.1};
\draw[densely dotted,black!36] (34.880,45.360) -- (35.880,45.360);
\node[anchor=west,inner sep=0pt,font=\rmfamily\fontsize{7}{7}\selectfont] at (36.880,45.360) {EF380018.1};
\draw[densely dotted,black!36] (34.880,48.600) -- (35.880,48.600);
\node[anchor=west,inner sep=0pt,font=\rmfamily\fontsize{7}{7}\selectfont] at (36.880,48.600) {GQ153914.1};
\draw[densely dotted,black!36] (34.880,51.840) -- (35.880,51.840);
\node[anchor=west,inner sep=0pt,font=\rmfamily\fontsize{7}{7}\selectfont] at (36.880,51.840) {GQ153913.1};
\draw[densely dotted,black!36] (34.880,55.080) -- (35.880,55.080);
\node[anchor=west,inner sep=0pt,font=\rmfamily\fontsize{7}{7}\selectfont] at (36.880,55.080) {EF380010.1};
\draw[densely dotted,black!36] (34.880,58.320) -- (35.880,58.320);
\node[anchor=west,inner sep=0pt,font=\rmfamily\fontsize{7}{7}\selectfont] at (36.880,58.320) {EF380016.1};
\draw[densely dotted,black!36] (34.880,61.560) -- (35.880,61.560);
\node[anchor=west,inner sep=0pt,font=\rmfamily\fontsize{7}{7}\selectfont] at (36.880,61.560) {EF380019.1};
\draw[densely dotted,black!36] (34.880,64.800) -- (35.880,64.800);
\node[anchor=west,inner sep=0pt,font=\rmfamily\fontsize{7}{7}\selectfont] at (36.880,64.800) {EF380020.1};
\draw[densely dotted,black!36] (34.880,68.040) -- (35.880,68.040);
\node[anchor=west,inner sep=0pt,font=\rmfamily\fontsize{7}{7}\selectfont] at (36.880,68.040) {EF380015.1};
\draw[densely dotted,black!36] (34.880,71.280) -- (35.880,71.280);
\node[anchor=west,inner sep=0pt,font=\rmfamily\fontsize{7}{7}\selectfont] at (36.880,71.280) {EF380017.1};
\draw[densely dotted,black!36] (34.880,74.520) -- (35.880,74.520);
\node[anchor=west,inner sep=0pt,font=\rmfamily\fontsize{7}{7}\selectfont] at (36.880,74.520) {GQ153912.1};
\draw[densely dotted,black!36] (34.880,77.760) -- (35.880,77.760);
\node[anchor=west,inner sep=0pt,font=\rmfamily\fontsize{7}{7}\selectfont] at (36.880,77.760) {AF176027.1};
\draw[densely dotted,black!36] (34.880,81.000) -- (35.880,81.000);
\node[anchor=west,inner sep=0pt,font=\rmfamily\fontsize{7}{7}\selectfont] at (36.880,81.000) {AF176033.1};
\draw[densely dotted,black!36] (34.880,84.240) -- (35.880,84.240);
\node[anchor=west,inner sep=0pt,font=\rmfamily\fontsize{7}{7}\selectfont] at (36.880,84.240) {AF176029.1};
\draw[densely dotted,black!36] (34.880,87.480) -- (35.880,87.480);
\node[anchor=west,inner sep=0pt,font=\rmfamily\fontsize{7}{7}\selectfont] at (36.880,87.480) {AF176031.1};
\draw[densely dotted,black!36] (34.880,90.720) -- (35.880,90.720);
\node[anchor=west,inner sep=0pt,font=\rmfamily\fontsize{7}{7}\selectfont] at (36.880,90.720) {AF299313.1};
\draw[densely dotted,black!36] (34.880,93.960) -- (35.880,93.960);
\node[anchor=west,inner sep=0pt,font=\rmfamily\fontsize{7}{7}\selectfont] at (36.880,93.960) {AF299301.1};
\draw[densely dotted,black!36] (34.880,97.200) -- (35.880,97.200);
\node[anchor=west,inner sep=0pt,font=\rmfamily\fontsize{7}{7}\selectfont] at (36.880,97.200) {AF299305.1};
\draw[densely dotted,black!36] (34.880,100.440) -- (35.880,100.440);
\node[anchor=west,inner sep=0pt,font=\rmfamily\fontsize{7}{7}\selectfont] at (36.880,100.440) {EF380009.1};
\draw[densely dotted,black!36] (34.880,103.680) -- (35.880,103.680);
\node[anchor=west,inner sep=0pt,font=\rmfamily\fontsize{7}{7}\selectfont] at (36.880,103.680) {EF380031.1};
\draw[densely dotted,black!36] (34.880,106.920) -- (35.880,106.920);
\node[anchor=west,inner sep=0pt,font=\rmfamily\fontsize{7}{7}\selectfont] at (36.880,106.920) {EF380032.1};
\draw[densely dotted,black!36] (34.880,110.160) -- (35.880,110.160);
\node[anchor=west,inner sep=0pt,font=\rmfamily\fontsize{7}{7}\selectfont] at (36.880,110.160) {GQ153915.1};
\draw[densely dotted,black!36] (34.880,113.400) -- (35.880,113.400);
\node[anchor=west,inner sep=0pt,font=\rmfamily\fontsize{7}{7}\selectfont] at (36.880,113.400) {NC\_001422.1};
\draw[densely dotted,black!36] (34.880,116.640) -- (35.880,116.640);
\node[anchor=west,inner sep=0pt,font=\rmfamily\fontsize{7}{7}\selectfont] at (36.880,116.640) {CP004084.1};
\draw[densely dotted,black!36] (34.880,119.880) -- (35.880,119.880);
\node[anchor=west,inner sep=0pt,font=\rmfamily\fontsize{7}{7}\selectfont] at (36.880,119.880) {AF299302.1};
\draw[densely dotted,black!36] (34.880,123.120) -- (35.880,123.120);
\node[anchor=west,inner sep=0pt,font=\rmfamily\fontsize{7}{7}\selectfont] at (36.880,123.120) {HM753662.1};
\draw[densely dotted,black!36] (34.880,126.360) -- (35.880,126.360);
\node[anchor=west,inner sep=0pt,font=\rmfamily\fontsize{7}{7}\selectfont] at (36.880,126.360) {GU385905.1};
\draw[densely dotted,black!36] (34.880,129.600) -- (35.880,129.600);
\node[anchor=west,inner sep=0pt,font=\rmfamily\fontsize{7}{7}\selectfont] at (36.880,129.600) {AF176034.1};
\draw[densely dotted,black!36] (34.880,132.840) -- (35.880,132.840);
\node[anchor=west,inner sep=0pt,font=\rmfamily\fontsize{7}{7}\selectfont] at (36.880,132.840) {FJ849058.1};
\draw[densely dotted,black!36] (34.880,136.080) -- (35.880,136.080);
\node[anchor=west,inner sep=0pt,font=\rmfamily\fontsize{7}{7}\selectfont] at (36.880,136.080) {HM753664.1};
\draw[densely dotted,black!36] (34.880,139.320) -- (35.880,139.320);
\node[anchor=west,inner sep=0pt,font=\rmfamily\fontsize{7}{7}\selectfont] at (36.880,139.320) {HM753665.1};
\draw[densely dotted,black!36] (34.880,142.560) -- (35.880,142.560);
\node[anchor=west,inner sep=0pt,font=\rmfamily\fontsize{7}{7}\selectfont] at (36.880,142.560) {HM753658.1};
\draw[densely dotted,black!36] (34.880,145.800) -- (35.880,145.800);
\node[anchor=west,inner sep=0pt,font=\rmfamily\fontsize{7}{7}\selectfont] at (36.880,145.800) {HM753660.1};
\draw[densely dotted,black!36] (34.880,149.040) -- (35.880,149.040);
\node[anchor=west,inner sep=0pt,font=\rmfamily\fontsize{7}{7}\selectfont] at (36.880,149.040) {HM753661.1};
\draw[densely dotted,black!36] (34.880,152.280) -- (35.880,152.280);
\node[anchor=west,inner sep=0pt,font=\rmfamily\fontsize{7}{7}\selectfont] at (36.880,152.280) {HM753666.1};
\draw[densely dotted,black!36] (34.880,155.520) -- (35.880,155.520);
\node[anchor=west,inner sep=0pt,font=\rmfamily\fontsize{7}{7}\selectfont] at (36.880,155.520) {HM753659.1};
\draw[densely dotted,black!36] (34.880,158.760) -- (35.880,158.760);
\node[anchor=west,inner sep=0pt,font=\rmfamily\fontsize{7}{7}\selectfont] at (36.880,158.760) {HM753663.1};
\draw[densely dotted,black!36] (34.880,162.000) -- (35.880,162.000);
\node[anchor=west,inner sep=0pt,font=\rmfamily\fontsize{7}{7}\selectfont] at (36.880,162.000) {AF299307.1};
\draw[densely dotted,black!36] (34.880,165.240) -- (35.880,165.240);
\node[anchor=west,inner sep=0pt,font=\rmfamily\fontsize{7}{7}\selectfont] at (36.880,165.240) {AF299306.1};
\draw[densely dotted,black!36] (34.880,168.480) -- (35.880,168.480);
\node[anchor=west,inner sep=0pt,font=\rmfamily\fontsize{7}{7}\selectfont] at (36.880,168.480) {AF299308.1};
\draw[densely dotted,black!36] (34.880,171.720) -- (35.880,171.720);
\node[anchor=west,inner sep=0pt,font=\rmfamily\fontsize{7}{7}\selectfont] at (36.880,171.720) {AF299312.1};
\draw[densely dotted,black!36] (34.880,174.960) -- (35.880,174.960);
\node[anchor=west,inner sep=0pt,font=\rmfamily\fontsize{7}{7}\selectfont] at (36.880,174.960) {AF299303.1};
\draw[densely dotted,black!36] (34.880,178.200) -- (35.880,178.200);
\node[anchor=west,inner sep=0pt,font=\rmfamily\fontsize{7}{7}\selectfont] at (36.880,178.200) {AF299311.1};
\draw[densely dotted,black!36] (34.880,181.440) -- (35.880,181.440);
\node[anchor=west,inner sep=0pt,font=\rmfamily\fontsize{7}{7}\selectfont] at (36.880,181.440) {AF176032.1};
\draw[densely dotted,black!36] (34.880,184.680) -- (35.880,184.680);
\node[anchor=west,inner sep=0pt,font=\rmfamily\fontsize{7}{7}\selectfont] at (36.880,184.680) {AF176028.1};
\draw[densely dotted,black!36] (34.880,187.920) -- (35.880,187.920);
\node[anchor=west,inner sep=0pt,font=\rmfamily\fontsize{7}{7}\selectfont] at (36.880,187.920) {AF176030.1};
\draw[densely dotted,black!36] (34.880,191.160) -- (35.880,191.160);
\node[anchor=west,inner sep=0pt,font=\rmfamily\fontsize{7}{7}\selectfont] at (36.880,191.160) {AF299314.1};
\draw[densely dotted,black!36] (34.880,194.400) -- (35.880,194.400);
\node[anchor=west,inner sep=0pt,font=\rmfamily\fontsize{7}{7}\selectfont] at (36.880,194.400) {AF299309.1};
\draw[densely dotted,black!36] (34.880,197.640) -- (35.880,197.640);
\node[anchor=west,inner sep=0pt,font=\rmfamily\fontsize{7}{7}\selectfont] at (36.880,197.640) {AF299310.1};
\draw[densely dotted,black!36] (34.880,200.880) -- (35.880,200.880);
\node[anchor=west,inner sep=0pt,font=\rmfamily\fontsize{7}{7}\selectfont] at (36.880,200.880) {AF299300.1};
\draw[densely dotted,black!36] (34.880,204.120) -- (35.880,204.120);
\node[anchor=west,inner sep=0pt,font=\rmfamily\fontsize{7}{7}\selectfont] at (36.880,204.120) {AF299304.1};
\draw[line width=.6pt] (0.000,210) -- ++(9.045,0) node[midway,below=2pt,font=\scriptsize] {10 éditions};
\node[anchor=west,font=\small\bfseries,inner sep=0pt] at (62.880,-6) {b) Neighbor-Joining};
\draw[draw=powderTeal!50!black] (62.887,0.000) -- (62.887,0.000);
\draw[draw=powderTeal!50!black] (62.887,3.240) -- (63.791,3.240);
\draw[draw=powderTeal!50!black] (62.887,0.000) -- (62.887,3.240);
\draw[draw=powderTeal!50!black] (62.880,1.620) -- (62.887,1.620);
\draw[draw=powderTeal!50!black] (63.777,6.480) -- (63.792,6.480);
\draw[draw=powderTeal!50!black] (64.666,9.720) -- (65.571,9.720);
\draw[draw=powderTeal!50!black] (64.666,12.960) -- (65.571,12.960);
\draw[draw=powderTeal!50!black] (64.666,9.720) -- (64.666,12.960);
\draw[draw=powderTeal!50!black] (63.777,11.340) -- (64.666,11.340);
\draw[draw=powderTeal!50!black] (63.777,6.480) -- (63.777,11.340);
\draw[draw=powderTeal!50!black] (62.887,8.910) -- (63.777,8.910);
\draw[draw=powderTeal!50!black] (62.904,16.200) -- (64.691,16.200);
\draw[draw=powderTeal!50!black] (64.603,19.440) -- (65.372,19.440);
\draw[draw=powderTeal!50!black] (75.779,22.680) -- (78.780,22.680);
\draw[draw=powderTeal!50!black] (84.536,25.920) -- (85.554,25.920);
\draw[draw=powderTeal!50!black] (85.327,29.160) -- (86.952,29.160);
\draw[draw=powderTeal!50!black] (85.327,32.400) -- (88.225,32.400);
\draw[draw=powderTeal!50!black] (85.327,29.160) -- (85.327,32.400);
\draw[draw=powderTeal!50!black] (84.536,30.780) -- (85.327,30.780);
\draw[draw=powderTeal!50!black] (84.536,25.920) -- (84.536,30.780);
\draw[draw=powderTeal!50!black] (75.779,28.350) -- (84.536,28.350);
\draw[draw=powderTeal!50!black] (75.779,22.680) -- (75.779,28.350);
\draw[draw=powderTeal!50!black] (69.377,25.515) -- (75.779,25.515);
\draw[draw=powderTeal!50!black] (113.436,35.640) -- (114.751,35.640);
\draw[draw=powderTeal!50!black] (114.835,38.880) -- (117.586,38.880);
\draw[draw=powderTeal!50!black] (114.835,42.120) -- (119.320,42.120);
\draw[draw=powderTeal!50!black] (114.835,38.880) -- (114.835,42.120);
\draw[draw=powderTeal!50!black] (113.436,40.500) -- (114.835,40.500);
\draw[draw=powderTeal!50!black] (113.436,35.640) -- (113.436,40.500);
\draw[draw=powderTeal!50!black] (92.589,38.070) -- (113.436,38.070);
\draw[draw=powderTeal!50!black] (95.894,45.360) -- (109.861,45.360);
\draw[draw=powderTeal!50!black] (106.473,48.600) -- (130.888,48.600);
\draw[draw=powderTeal!50!black] (119.142,51.840) -- (120.402,51.840);
\draw[draw=powderTeal!50!black] (120.595,55.080) -- (125.992,55.080);
\draw[draw=powderTeal!50!black] (120.595,58.320) -- (136.000,58.320);
\draw[draw=powderTeal!50!black] (120.595,55.080) -- (120.595,58.320);
\draw[draw=powderTeal!50!black] (119.142,56.700) -- (120.595,56.700);
\draw[draw=powderTeal!50!black] (119.142,51.840) -- (119.142,56.700);
\draw[draw=powderTeal!50!black] (106.473,54.270) -- (119.142,54.270);
\draw[draw=powderTeal!50!black] (106.473,48.600) -- (106.473,54.270);
\draw[draw=powderTeal!50!black] (96.681,51.435) -- (106.473,51.435);
\draw[draw=powderTeal!50!black] (99.213,61.560) -- (100.225,61.560);
\draw[draw=powderTeal!50!black] (104.532,64.800) -- (104.768,64.800);
\draw[draw=powderTeal!50!black] (105.200,68.040) -- (112.521,68.040);
\draw[draw=powderTeal!50!black] (105.200,71.280) -- (111.447,71.280);
\draw[draw=powderTeal!50!black] (105.200,68.040) -- (105.200,71.280);
\draw[draw=powderTeal!50!black] (104.532,69.660) -- (105.200,69.660);
\draw[draw=powderTeal!50!black] (104.532,64.800) -- (104.532,69.660);
\draw[draw=powderTeal!50!black] (99.213,67.230) -- (104.532,67.230);
\draw[draw=powderTeal!50!black] (99.213,61.560) -- (99.213,67.230);
\draw[draw=powderTeal!50!black] (96.681,64.395) -- (99.213,64.395);
\draw[draw=powderTeal!50!black] (96.681,51.435) -- (96.681,64.395);
\draw[draw=powderTeal!50!black] (95.894,57.915) -- (96.681,57.915);
\draw[draw=powderTeal!50!black] (95.894,45.360) -- (95.894,57.915);
\draw[draw=powderTeal!50!black] (92.589,51.638) -- (95.894,51.638);
\draw[draw=powderTeal!50!black] (92.589,38.070) -- (92.589,51.638);
\draw[draw=powderTeal!50!black] (71.782,44.854) -- (92.589,44.854);
\draw[draw=powderTeal!50!black] (81.948,74.520) -- (91.095,74.520);
\draw[draw=powderTeal!50!black] (104.004,77.760) -- (104.106,77.760);
\draw[draw=powderTeal!50!black] (104.004,81.000) -- (110.234,81.000);
\draw[draw=powderTeal!50!black] (104.004,77.760) -- (104.004,81.000);
\draw[draw=powderTeal!50!black] (81.948,79.380) -- (104.004,79.380);
\draw[draw=powderTeal!50!black] (81.948,74.520) -- (81.948,79.380);
\draw[draw=powderTeal!50!black] (77.017,76.950) -- (81.948,76.950);
\draw[draw=powderTeal!50!black] (78.079,84.240) -- (78.155,84.240);
\draw[draw=powderTeal!50!black] (78.079,87.480) -- (81.620,87.480);
\draw[draw=powderTeal!50!black] (78.079,84.240) -- (78.079,87.480);
\draw[draw=powderTeal!50!black] (77.017,85.860) -- (78.079,85.860);
\draw[draw=powderTeal!50!black] (77.017,76.950) -- (77.017,85.860);
\draw[draw=powderTeal!50!black] (71.782,81.405) -- (77.017,81.405);
\draw[draw=powderTeal!50!black] (71.782,44.854) -- (71.782,81.405);
\draw[draw=powderTeal!50!black] (69.377,63.129) -- (71.782,63.129);
\draw[draw=powderTeal!50!black] (69.377,25.515) -- (69.377,63.129);
\draw[draw=powderTeal!50!black] (67.071,44.322) -- (69.377,44.322);
\draw[draw=powderTeal!50!black] (67.390,90.720) -- (69.264,90.720);
\draw[draw=powderTeal!50!black] (73.658,93.960) -- (73.752,93.960);
\draw[draw=powderTeal!50!black] (74.242,97.200) -- (95.196,97.200);
\draw[draw=powderTeal!50!black] (74.995,100.440) -- (87.119,100.440);
\draw[draw=powderTeal!50!black] (74.995,103.680) -- (95.432,103.680);
\draw[draw=powderTeal!50!black] (74.995,100.440) -- (74.995,103.680);
\draw[draw=powderTeal!50!black] (74.242,102.060) -- (74.995,102.060);
\draw[draw=powderTeal!50!black] (74.242,97.200) -- (74.242,102.060);
\draw[draw=powderTeal!50!black] (73.658,99.630) -- (74.242,99.630);
\draw[draw=powderTeal!50!black] (73.658,93.960) -- (73.658,99.630);
\draw[draw=powderTeal!50!black] (67.390,96.795) -- (73.658,96.795);
\draw[draw=powderTeal!50!black] (67.390,90.720) -- (67.390,96.795);
\draw[draw=powderTeal!50!black] (67.071,93.757) -- (67.390,93.757);
\draw[draw=powderTeal!50!black] (67.071,44.322) -- (67.071,93.757);
\draw[draw=powderTeal!50!black] (65.434,69.040) -- (67.071,69.040);
\draw[draw=powderTeal!50!black] (65.682,106.920) -- (65.845,106.920);
\draw[draw=powderTeal!50!black] (66.007,110.160) -- (67.078,110.160);
\draw[draw=powderTeal!50!black] (66.617,113.400) -- (67.541,113.400);
\draw[draw=powderTeal!50!black] (72.290,116.640) -- (73.039,116.640);
\draw[draw=powderTeal!50!black] (89.214,119.880) -- (105.073,119.880);
\draw[draw=powderTeal!50!black] (89.214,123.120) -- (111.343,123.120);
\draw[draw=powderTeal!50!black] (89.214,119.880) -- (89.214,123.120);
\draw[draw=powderTeal!50!black] (73.589,121.500) -- (89.214,121.500);
\draw[draw=powderTeal!50!black] (76.149,126.360) -- (96.536,126.360);
\draw[draw=powderTeal!50!black] (78.691,129.600) -- (96.047,129.600);
\draw[draw=powderTeal!50!black] (92.088,132.840) -- (94.326,132.840);
\draw[draw=powderTeal!50!black] (99.798,136.080) -- (118.515,136.080);
\draw[draw=powderTeal!50!black] (101.885,139.320) -- (110.257,139.320);
\draw[draw=powderTeal!50!black] (101.885,142.560) -- (110.698,142.560);
\draw[draw=powderTeal!50!black] (101.885,139.320) -- (101.885,142.560);
\draw[draw=powderTeal!50!black] (99.798,140.940) -- (101.885,140.940);
\draw[draw=powderTeal!50!black] (99.798,136.080) -- (99.798,140.940);
\draw[draw=powderTeal!50!black] (92.088,138.510) -- (99.798,138.510);
\draw[draw=powderTeal!50!black] (92.088,132.840) -- (92.088,138.510);
\draw[draw=powderTeal!50!black] (78.691,135.675) -- (92.088,135.675);
\draw[draw=powderTeal!50!black] (78.691,129.600) -- (78.691,135.675);
\draw[draw=powderTeal!50!black] (77.525,132.638) -- (78.691,132.638);
\draw[draw=powderTeal!50!black] (78.055,145.800) -- (100.571,145.800);
\draw[draw=powderTeal!50!black] (78.055,149.040) -- (100.762,149.040);
\draw[draw=powderTeal!50!black] (78.055,145.800) -- (78.055,149.040);
\draw[draw=powderTeal!50!black] (77.525,147.420) -- (78.055,147.420);
\draw[draw=powderTeal!50!black] (77.525,132.638) -- (77.525,147.420);
\draw[draw=powderTeal!50!black] (76.149,140.029) -- (77.525,140.029);
\draw[draw=powderTeal!50!black] (76.149,126.360) -- (76.149,140.029);
\draw[draw=powderTeal!50!black] (75.229,133.194) -- (76.149,133.194);
\draw[draw=powderTeal!50!black] (75.246,152.280) -- (94.723,152.280);
\draw[draw=powderTeal!50!black] (76.566,155.520) -- (96.476,155.520);
\draw[draw=powderTeal!50!black] (76.566,158.760) -- (96.453,158.760);
\draw[draw=powderTeal!50!black] (76.566,155.520) -- (76.566,158.760);
\draw[draw=powderTeal!50!black] (75.493,157.140) -- (76.566,157.140);
\draw[draw=powderTeal!50!black] (77.460,162.000) -- (102.542,162.000);
\draw[draw=powderTeal!50!black] (77.460,165.240) -- (92.175,165.240);
\draw[draw=powderTeal!50!black] (77.460,162.000) -- (77.460,165.240);
\draw[draw=powderTeal!50!black] (75.911,163.620) -- (77.460,163.620);
\draw[draw=powderTeal!50!black] (80.957,168.480) -- (113.848,168.480);
\draw[draw=powderTeal!50!black] (81.530,171.720) -- (112.307,171.720);
\draw[draw=powderTeal!50!black] (82.410,174.960) -- (111.363,174.960);
\draw[draw=powderTeal!50!black] (82.410,178.200) -- (108.630,178.200);
\draw[draw=powderTeal!50!black] (82.410,174.960) -- (82.410,178.200);
\draw[draw=powderTeal!50!black] (81.530,176.580) -- (82.410,176.580);
\draw[draw=powderTeal!50!black] (81.530,171.720) -- (81.530,176.580);
\draw[draw=powderTeal!50!black] (80.957,174.150) -- (81.530,174.150);
\draw[draw=powderTeal!50!black] (80.957,168.480) -- (80.957,174.150);
\draw[draw=powderTeal!50!black] (79.830,171.315) -- (80.957,171.315);
\draw[draw=powderTeal!50!black] (80.513,181.440) -- (82.312,181.440);
\draw[draw=powderTeal!50!black] (80.513,184.680) -- (81.427,184.680);
\draw[draw=powderTeal!50!black] (80.513,181.440) -- (80.513,184.680);
\draw[draw=powderTeal!50!black] (79.830,183.060) -- (80.513,183.060);
\draw[draw=powderTeal!50!black] (79.830,171.315) -- (79.830,183.060);
\draw[draw=powderTeal!50!black] (76.963,177.188) -- (79.830,177.188);
\draw[draw=powderTeal!50!black] (78.221,187.920) -- (102.157,187.920);
\draw[draw=powderTeal!50!black] (78.221,191.160) -- (107.649,191.160);
\draw[draw=powderTeal!50!black] (78.221,187.920) -- (78.221,191.160);
\draw[draw=powderTeal!50!black] (76.963,189.540) -- (78.221,189.540);
\draw[draw=powderTeal!50!black] (76.963,177.188) -- (76.963,189.540);
\draw[draw=powderTeal!50!black] (75.911,183.364) -- (76.963,183.364);
\draw[draw=powderTeal!50!black] (75.911,163.620) -- (75.911,183.364);
\draw[draw=powderTeal!50!black] (75.493,173.492) -- (75.911,173.492);
\draw[draw=powderTeal!50!black] (75.493,157.140) -- (75.493,173.492);
\draw[draw=powderTeal!50!black] (75.246,165.316) -- (75.493,165.316);
\draw[draw=powderTeal!50!black] (75.246,152.280) -- (75.246,165.316);
\draw[draw=powderTeal!50!black] (75.229,158.798) -- (75.246,158.798);
\draw[draw=powderTeal!50!black] (75.229,133.194) -- (75.229,158.798);
\draw[draw=powderTeal!50!black] (73.589,145.996) -- (75.229,145.996);
\draw[draw=powderTeal!50!black] (73.589,121.500) -- (73.589,145.996);
\draw[draw=powderTeal!50!black] (72.290,133.748) -- (73.589,133.748);
\draw[draw=powderTeal!50!black] (72.290,116.640) -- (72.290,133.748);
\draw[draw=powderTeal!50!black] (67.050,125.194) -- (72.290,125.194);
\draw[draw=powderTeal!50!black] (68.593,194.400) -- (70.459,194.400);
\draw[draw=powderTeal!50!black] (68.593,197.640) -- (71.250,197.640);
\draw[draw=powderTeal!50!black] (68.593,194.400) -- (68.593,197.640);
\draw[draw=powderTeal!50!black] (67.050,196.020) -- (68.593,196.020);
\draw[draw=powderTeal!50!black] (67.050,125.194) -- (67.050,196.020);
\draw[draw=powderTeal!50!black] (66.617,160.607) -- (67.050,160.607);
\draw[draw=powderTeal!50!black] (66.617,113.400) -- (66.617,160.607);
\draw[draw=powderTeal!50!black] (66.007,137.004) -- (66.617,137.004);
\draw[draw=powderTeal!50!black] (66.007,110.160) -- (66.007,137.004);
\draw[draw=powderTeal!50!black] (65.682,123.582) -- (66.007,123.582);
\draw[draw=powderTeal!50!black] (65.682,106.920) -- (65.682,123.582);
\draw[draw=powderTeal!50!black] (65.434,115.251) -- (65.682,115.251);
\draw[draw=powderTeal!50!black] (65.434,69.040) -- (65.434,115.251);
\draw[draw=powderTeal!50!black] (64.603,92.145) -- (65.434,92.145);
\draw[draw=powderTeal!50!black] (64.603,19.440) -- (64.603,92.145);
\draw[draw=powderTeal!50!black] (63.373,55.793) -- (64.603,55.793);
\draw[draw=powderTeal!50!black] (63.798,200.880) -- (63.798,200.880);
\draw[draw=powderTeal!50!black] (63.798,204.120) -- (63.798,204.120);
\draw[draw=powderTeal!50!black] (63.798,200.880) -- (63.798,204.120);
\draw[draw=powderTeal!50!black] (63.373,202.500) -- (63.798,202.500);
\draw[draw=powderTeal!50!black] (63.373,55.793) -- (63.373,202.500);
\draw[draw=powderTeal!50!black] (62.904,129.146) -- (63.373,129.146);
\draw[draw=powderTeal!50!black] (62.904,16.200) -- (62.904,129.146);
\draw[draw=powderTeal!50!black] (62.887,72.673) -- (62.904,72.673);
\draw[draw=powderTeal!50!black] (62.887,8.910) -- (62.887,72.673);
\draw[draw=powderTeal!50!black] (62.880,40.792) -- (62.887,40.792);
\draw[draw=powderTeal!50!black] (62.880,1.620) -- (62.880,40.792);
\draw[densely dotted,black!36] (62.887,0.000) -- (137.000,0.000);
\node[anchor=west,inner sep=0pt,font=\rmfamily\fontsize{7}{7}\selectfont] at (138.000,0.000) {HM753659.1};
\draw[densely dotted,black!36] (63.791,3.240) -- (137.000,3.240);
\node[anchor=west,inner sep=0pt,font=\rmfamily\fontsize{7}{7}\selectfont] at (138.000,3.240) {HM753663.1};
\draw[densely dotted,black!36] (63.792,6.480) -- (137.000,6.480);
\node[anchor=west,inner sep=0pt,font=\rmfamily\fontsize{7}{7}\selectfont] at (138.000,6.480) {HM753666.1};
\draw[densely dotted,black!36] (65.571,9.720) -- (137.000,9.720);
\node[anchor=west,inner sep=0pt,font=\rmfamily\fontsize{7}{7}\selectfont] at (138.000,9.720) {HM753664.1};
\draw[densely dotted,black!36] (65.571,12.960) -- (137.000,12.960);
\node[anchor=west,inner sep=0pt,font=\rmfamily\fontsize{7}{7}\selectfont] at (138.000,12.960) {HM753665.1};
\draw[densely dotted,black!36] (64.691,16.200) -- (137.000,16.200);
\node[anchor=west,inner sep=0pt,font=\rmfamily\fontsize{7}{7}\selectfont] at (138.000,16.200) {HM753658.1};
\draw[densely dotted,black!36] (65.372,19.440) -- (137.000,19.440);
\node[anchor=west,inner sep=0pt,font=\rmfamily\fontsize{7}{7}\selectfont] at (138.000,19.440) {HM753662.1};
\draw[densely dotted,black!36] (78.780,22.680) -- (137.000,22.680);
\node[anchor=west,inner sep=0pt,font=\rmfamily\fontsize{7}{7}\selectfont] at (138.000,22.680) {AF176027.1};
\draw[densely dotted,black!36] (85.554,25.920) -- (137.000,25.920);
\node[anchor=west,inner sep=0pt,font=\rmfamily\fontsize{7}{7}\selectfont] at (138.000,25.920) {AF176029.1};
\draw[densely dotted,black!36] (86.952,29.160) -- (137.000,29.160);
\node[anchor=west,inner sep=0pt,font=\rmfamily\fontsize{7}{7}\selectfont] at (138.000,29.160) {AF176031.1};
\draw[densely dotted,black!36] (88.225,32.400) -- (137.000,32.400);
\node[anchor=west,inner sep=0pt,font=\rmfamily\fontsize{7}{7}\selectfont] at (138.000,32.400) {AF176033.1};
\draw[densely dotted,black!36] (114.751,35.640) -- (137.000,35.640);
\node[anchor=west,inner sep=0pt,font=\rmfamily\fontsize{7}{7}\selectfont] at (138.000,35.640) {AF176028.1};
\draw[densely dotted,black!36] (117.586,38.880) -- (137.000,38.880);
\node[anchor=west,inner sep=0pt,font=\rmfamily\fontsize{7}{7}\selectfont] at (138.000,38.880) {AF176030.1};
\draw[densely dotted,black!36] (119.320,42.120) -- (137.000,42.120);
\node[anchor=west,inner sep=0pt,font=\rmfamily\fontsize{7}{7}\selectfont] at (138.000,42.120) {AF176032.1};
\draw[densely dotted,black!36] (109.861,45.360) -- (137.000,45.360);
\node[anchor=west,inner sep=0pt,font=\rmfamily\fontsize{7}{7}\selectfont] at (138.000,45.360) {AF299314.1};
\draw[densely dotted,black!36] (130.888,48.600) -- (137.000,48.600);
\node[anchor=west,inner sep=0pt,font=\rmfamily\fontsize{7}{7}\selectfont] at (138.000,48.600) {AF299307.1};
\draw[densely dotted,black!36] (120.402,51.840) -- (137.000,51.840);
\node[anchor=west,inner sep=0pt,font=\rmfamily\fontsize{7}{7}\selectfont] at (138.000,51.840) {AF299303.1};
\draw[densely dotted,black!36] (125.992,55.080) -- (137.000,55.080);
\node[anchor=west,inner sep=0pt,font=\rmfamily\fontsize{7}{7}\selectfont] at (138.000,55.080) {AF299311.1};
\draw[densely dotted,black!36] (136.000,58.320) -- (137.000,58.320);
\node[anchor=west,inner sep=0pt,font=\rmfamily\fontsize{7}{7}\selectfont] at (138.000,58.320) {AF299312.1};
\draw[densely dotted,black!36] (100.225,61.560) -- (137.000,61.560);
\node[anchor=west,inner sep=0pt,font=\rmfamily\fontsize{7}{7}\selectfont] at (138.000,61.560) {AF299300.1};
\draw[densely dotted,black!36] (104.768,64.800) -- (137.000,64.800);
\node[anchor=west,inner sep=0pt,font=\rmfamily\fontsize{7}{7}\selectfont] at (138.000,64.800) {AF299304.1};
\draw[densely dotted,black!36] (112.521,68.040) -- (137.000,68.040);
\node[anchor=west,inner sep=0pt,font=\rmfamily\fontsize{7}{7}\selectfont] at (138.000,68.040) {AF299309.1};
\draw[densely dotted,black!36] (111.447,71.280) -- (137.000,71.280);
\node[anchor=west,inner sep=0pt,font=\rmfamily\fontsize{7}{7}\selectfont] at (138.000,71.280) {AF299310.1};
\draw[densely dotted,black!36] (91.095,74.520) -- (137.000,74.520);
\node[anchor=west,inner sep=0pt,font=\rmfamily\fontsize{7}{7}\selectfont] at (138.000,74.520) {AF299313.1};
\draw[densely dotted,black!36] (104.106,77.760) -- (137.000,77.760);
\node[anchor=west,inner sep=0pt,font=\rmfamily\fontsize{7}{7}\selectfont] at (138.000,77.760) {AF299306.1};
\draw[densely dotted,black!36] (110.234,81.000) -- (137.000,81.000);
\node[anchor=west,inner sep=0pt,font=\rmfamily\fontsize{7}{7}\selectfont] at (138.000,81.000) {AF299308.1};
\draw[densely dotted,black!36] (78.155,84.240) -- (137.000,84.240);
\node[anchor=west,inner sep=0pt,font=\rmfamily\fontsize{7}{7}\selectfont] at (138.000,84.240) {AF299301.1};
\draw[densely dotted,black!36] (81.620,87.480) -- (137.000,87.480);
\node[anchor=west,inner sep=0pt,font=\rmfamily\fontsize{7}{7}\selectfont] at (138.000,87.480) {AF299305.1};
\draw[densely dotted,black!36] (69.264,90.720) -- (137.000,90.720);
\node[anchor=west,inner sep=0pt,font=\rmfamily\fontsize{7}{7}\selectfont] at (138.000,90.720) {AF299302.1};
\draw[densely dotted,black!36] (73.752,93.960) -- (137.000,93.960);
\node[anchor=west,inner sep=0pt,font=\rmfamily\fontsize{7}{7}\selectfont] at (138.000,93.960) {GQ153915.1};
\draw[densely dotted,black!36] (95.196,97.200) -- (137.000,97.200);
\node[anchor=west,inner sep=0pt,font=\rmfamily\fontsize{7}{7}\selectfont] at (138.000,97.200) {GQ153913.1};
\draw[densely dotted,black!36] (87.119,100.440) -- (137.000,100.440);
\node[anchor=west,inner sep=0pt,font=\rmfamily\fontsize{7}{7}\selectfont] at (138.000,100.440) {GQ153912.1};
\draw[densely dotted,black!36] (95.432,103.680) -- (137.000,103.680);
\node[anchor=west,inner sep=0pt,font=\rmfamily\fontsize{7}{7}\selectfont] at (138.000,103.680) {GQ153914.1};
\draw[densely dotted,black!36] (65.845,106.920) -- (137.000,106.920);
\node[anchor=west,inner sep=0pt,font=\rmfamily\fontsize{7}{7}\selectfont] at (138.000,106.920) {AF176034.1};
\draw[densely dotted,black!36] (67.078,110.160) -- (137.000,110.160);
\node[anchor=west,inner sep=0pt,font=\rmfamily\fontsize{7}{7}\selectfont] at (138.000,110.160) {FJ849058.1};
\draw[densely dotted,black!36] (67.541,113.400) -- (137.000,113.400);
\node[anchor=west,inner sep=0pt,font=\rmfamily\fontsize{7}{7}\selectfont] at (138.000,113.400) {GU385905.1};
\draw[densely dotted,black!36] (73.039,116.640) -- (137.000,116.640);
\node[anchor=west,inner sep=0pt,font=\rmfamily\fontsize{7}{7}\selectfont] at (138.000,116.640) {EF380009.1};
\draw[densely dotted,black!36] (105.073,119.880) -- (137.000,119.880);
\node[anchor=west,inner sep=0pt,font=\rmfamily\fontsize{7}{7}\selectfont] at (138.000,119.880) {EF380024.1};
\draw[densely dotted,black!36] (111.343,123.120) -- (137.000,123.120);
\node[anchor=west,inner sep=0pt,font=\rmfamily\fontsize{7}{7}\selectfont] at (138.000,123.120) {EF380025.1};
\draw[densely dotted,black!36] (96.536,126.360) -- (137.000,126.360);
\node[anchor=west,inner sep=0pt,font=\rmfamily\fontsize{7}{7}\selectfont] at (138.000,126.360) {EF380010.1};
\draw[densely dotted,black!36] (96.047,129.600) -- (137.000,129.600);
\node[anchor=west,inner sep=0pt,font=\rmfamily\fontsize{7}{7}\selectfont] at (138.000,129.600) {EF380020.1};
\draw[densely dotted,black!36] (94.326,132.840) -- (137.000,132.840);
\node[anchor=west,inner sep=0pt,font=\rmfamily\fontsize{7}{7}\selectfont] at (138.000,132.840) {EF380021.1};
\draw[densely dotted,black!36] (118.515,136.080) -- (137.000,136.080);
\node[anchor=west,inner sep=0pt,font=\rmfamily\fontsize{7}{7}\selectfont] at (138.000,136.080) {EF380029.1};
\draw[densely dotted,black!36] (110.257,139.320) -- (137.000,139.320);
\node[anchor=west,inner sep=0pt,font=\rmfamily\fontsize{7}{7}\selectfont] at (138.000,139.320) {EF380028.1};
\draw[densely dotted,black!36] (110.698,142.560) -- (137.000,142.560);
\node[anchor=west,inner sep=0pt,font=\rmfamily\fontsize{7}{7}\selectfont] at (138.000,142.560) {EF380030.1};
\draw[densely dotted,black!36] (100.571,145.800) -- (137.000,145.800);
\node[anchor=west,inner sep=0pt,font=\rmfamily\fontsize{7}{7}\selectfont] at (138.000,145.800) {EF380014.1};
\draw[densely dotted,black!36] (100.762,149.040) -- (137.000,149.040);
\node[anchor=west,inner sep=0pt,font=\rmfamily\fontsize{7}{7}\selectfont] at (138.000,149.040) {EF380018.1};
\draw[densely dotted,black!36] (94.723,152.280) -- (137.000,152.280);
\node[anchor=west,inner sep=0pt,font=\rmfamily\fontsize{7}{7}\selectfont] at (138.000,152.280) {EF380015.1};
\draw[densely dotted,black!36] (96.476,155.520) -- (137.000,155.520);
\node[anchor=west,inner sep=0pt,font=\rmfamily\fontsize{7}{7}\selectfont] at (138.000,155.520) {EF380016.1};
\draw[densely dotted,black!36] (96.453,158.760) -- (137.000,158.760);
\node[anchor=west,inner sep=0pt,font=\rmfamily\fontsize{7}{7}\selectfont] at (138.000,158.760) {EF380019.1};
\draw[densely dotted,black!36] (102.542,162.000) -- (137.000,162.000);
\node[anchor=west,inner sep=0pt,font=\rmfamily\fontsize{7}{7}\selectfont] at (138.000,162.000) {EF380013.1};
\draw[densely dotted,black!36] (92.175,165.240) -- (137.000,165.240);
\node[anchor=west,inner sep=0pt,font=\rmfamily\fontsize{7}{7}\selectfont] at (138.000,165.240) {EF380017.1};
\draw[densely dotted,black!36] (113.848,168.480) -- (137.000,168.480);
\node[anchor=west,inner sep=0pt,font=\rmfamily\fontsize{7}{7}\selectfont] at (138.000,168.480) {EF380022.1};
\draw[densely dotted,black!36] (112.307,171.720) -- (137.000,171.720);
\node[anchor=west,inner sep=0pt,font=\rmfamily\fontsize{7}{7}\selectfont] at (138.000,171.720) {EF380026.1};
\draw[densely dotted,black!36] (111.363,174.960) -- (137.000,174.960);
\node[anchor=west,inner sep=0pt,font=\rmfamily\fontsize{7}{7}\selectfont] at (138.000,174.960) {EF380023.1};
\draw[densely dotted,black!36] (108.630,178.200) -- (137.000,178.200);
\node[anchor=west,inner sep=0pt,font=\rmfamily\fontsize{7}{7}\selectfont] at (138.000,178.200) {EF380027.1};
\draw[densely dotted,black!36] (82.312,181.440) -- (137.000,181.440);
\node[anchor=west,inner sep=0pt,font=\rmfamily\fontsize{7}{7}\selectfont] at (138.000,181.440) {EF380031.1};
\draw[densely dotted,black!36] (81.427,184.680) -- (137.000,184.680);
\node[anchor=west,inner sep=0pt,font=\rmfamily\fontsize{7}{7}\selectfont] at (138.000,184.680) {EF380032.1};
\draw[densely dotted,black!36] (102.157,187.920) -- (137.000,187.920);
\node[anchor=west,inner sep=0pt,font=\rmfamily\fontsize{7}{7}\selectfont] at (138.000,187.920) {EF380011.1};
\draw[densely dotted,black!36] (107.649,191.160) -- (137.000,191.160);
\node[anchor=west,inner sep=0pt,font=\rmfamily\fontsize{7}{7}\selectfont] at (138.000,191.160) {EF380012.1};
\draw[densely dotted,black!36] (70.459,194.400) -- (137.000,194.400);
\node[anchor=west,inner sep=0pt,font=\rmfamily\fontsize{7}{7}\selectfont] at (138.000,194.400) {NC\_001422.1};
\draw[densely dotted,black!36] (71.250,197.640) -- (137.000,197.640);
\node[anchor=west,inner sep=0pt,font=\rmfamily\fontsize{7}{7}\selectfont] at (138.000,197.640) {CP004084.1};
\draw[densely dotted,black!36] (63.798,200.880) -- (137.000,200.880);
\node[anchor=west,inner sep=0pt,font=\rmfamily\fontsize{7}{7}\selectfont] at (138.000,200.880) {HM753660.1};
\draw[densely dotted,black!36] (63.798,204.120) -- (137.000,204.120);
\node[anchor=west,inner sep=0pt,font=\rmfamily\fontsize{7}{7}\selectfont] at (138.000,204.120) {HM753661.1};
\draw[line width=.6pt] (62.880,210) -- ++(9.045,0) node[midway,below=2pt,font=\scriptsize] {10 éditions};
\end{tikzpicture}
}
\caption{Les deux phylogrammes portent les identifiants des séquences et une barre horizontale de dix éditions pour chaque arbre.}\label{fig:arbres}\sourcefig{report/phylograms.py}
\end{figure}
\synthese{L'écart absolu maximal vaut 56,87 éditions pour UPGMA et 63,86 pour NJ.}{L'erreur maximale et l'ajustement moyen de la figure~\ref{fig:phage} ordonnent donc les arbres différemment.}

\chapitre{Structure des groupes dans les arbres de surveillance}
Les distances entre 64 séquences archivées de SARS-CoV-2 permettent d'examiner si les groupes étiquetés sont séparés par une arête. Dans un arbre enraciné, un groupe est retrouvé si les descendants d'une arête positive coïncident avec ses membres. Pour NJ, non enraciné, les deux côtés de l'arête sont examinés. Sur les 1\,905\,372 caractères lus, 65\,588 sont ambigus. La distance compte les différences entre caractères déposés, y compris aux positions ambiguës.
\chart{19_phylogenie_surveillance.png}{Les distances globales entre paires et les groupes isolés par une arête sont comparés aux partitions de l'arbre de surveillance sur les identifiants communs.}{fig:surveillance}{scripts/phylo_groups.py}
Les 2\,016 scores de NvPD correspondent à ceux d'edlib et de WFA2. Sur M1 Pro, six processus NvPD monofils qui se partagent les paires terminent en 268,83\,s. La commande qui attribue huit fils à chaque paire atteint la limite de 1\,200\,s. À la figure~\ref{fig:surveillance}, la médiane est de 112 éditions pour les 252 paires de même lignage et de 402 pour les 1\,701 paires de lignages différents. Ces statistiques excluent les 63 paires qui contiennent la référence. Sur 52 feuilles communes, l'arbre public UShER~\cite{usher} donne des longueurs de chemin en événements portés par l'arbre. Leur corrélation descriptive avec les distances d'édition vaut \(r=0,124\) sur 1\,326 paires qui partagent des feuilles.
\synthese{Sur la cohorte entière, UPGMA isole deux des sept groupes et NJ en isole trois.}{Sur les 52 feuilles partagées, les arbres recalculés retrouvent respectivement six et huit des 27 partitions internes positives de l'arbre public.}
